\documentclass{standalone}
\usepackage{circuitikz}
\usetikzlibrary{calc}
\definecolor{circuitnotemark}{HTML}{FE00FE}
\begin{document}
\begin{circuitikz}[american, line width=0.8pt]
\ctikzset{bipoles/length=1.2cm}
\foreach \x in {0,1.2,2.4,3.6,4.8} {\foreach \y in {0,-1.2,-2.4,-3.6,-4.8,-6,-7.2,-8.4,-9.6,-10.8,-12} {\fill[gray, opacity=0.35] (\x,\y) circle (0.66pt);}}
\node[gray, font=\scriptsize] at (0,0.504) {1};
\node[gray, font=\scriptsize] at (1.2,0.504) {2};
\node[gray, font=\scriptsize] at (2.4,0.504) {3};
\node[gray, font=\scriptsize] at (3.6,0.504) {4};
\node[gray, font=\scriptsize] at (4.8,0.504) {5};
\node[gray, font=\scriptsize, anchor=east] at (-0.504,0) {a};
\node[gray, font=\scriptsize, anchor=east] at (-0.504,-1.2) {b};
\node[gray, font=\scriptsize, anchor=east] at (-0.504,-2.4) {c};
\node[gray, font=\scriptsize, anchor=east] at (-0.504,-3.6) {d};
\node[gray, font=\scriptsize, anchor=east] at (-0.504,-4.8) {e};
\node[gray, font=\scriptsize, anchor=east] at (-0.504,-6) {f};
\node[gray, font=\scriptsize, anchor=east] at (-0.504,-7.2) {g};
\node[gray, font=\scriptsize, anchor=east] at (-0.504,-8.4) {h};
\node[gray, font=\scriptsize, anchor=east] at (-0.504,-9.6) {i};
\node[gray, font=\scriptsize, anchor=east] at (-0.504,-10.8) {j};
\node[gray, font=\scriptsize, anchor=east] at (-0.504,-12) {k};
\coordinate (b2) at (1.2,-1.2);
\coordinate (b5) at (4.8,-1.2);
\coordinate (f2) at (1.2,-6);
\coordinate (f5) at (4.8,-6);
\coordinate (j2) at (1.2,-10.8);
\coordinate (j5) at (4.8,-10.8);
\coordinate (a2) at (1.2,0);
\coordinate (c2) at (1.2,-2.4);
\coordinate (b1) at (0,-1.2);
\coordinate (a5) at (4.8,0);
\coordinate (c5) at (4.8,-2.4);
\coordinate (b4) at (3.6,-1.2);
\coordinate (e2) at (1.2,-4.8);
\coordinate (g2) at (1.2,-7.2);
\coordinate (f1) at (0,-6);
\coordinate (e5) at (4.8,-4.8);
\coordinate (g5) at (4.8,-7.2);
\coordinate (f4) at (3.6,-6);
\coordinate (i2) at (1.2,-9.6);
\coordinate (k2) at (1.2,-12);
\coordinate (j1) at (0,-10.8);
\coordinate (i5) at (4.8,-9.6);
\coordinate (k5) at (4.8,-12);
\coordinate (j4) at (3.6,-10.8);
\node[njfet] (part-J1) at (b2) {}; % line 2
\node[pjfet] (part-J2) at (b5) {}; % line 3
\node[nigfete] (part-M1) at (f2) {}; % line 4
\node[pigfete] (part-M2) at (f5) {}; % line 5
\node[nigfetd] (part-M3) at (j2) {}; % line 6
\node[pigfetd] (part-M4) at (j5) {}; % line 7
\draw (a2) -| (part-J1.drain); % line 9
\draw (c2) -| (part-J1.source); % line 10
\draw (b1) -| (part-J1.gate); % line 11
\draw (a5) -| (part-J2.drain); % line 12
\draw (c5) -| (part-J2.source); % line 13
\draw (b4) -| (part-J2.gate); % line 14
\draw (e2) -| (part-M1.drain); % line 15
\draw (g2) -| (part-M1.source); % line 16
\draw (f1) -| (part-M1.gate); % line 17
\draw (e5) -| (part-M2.drain); % line 18
\draw (g5) -| (part-M2.source); % line 19
\draw (f4) -| (part-M2.gate); % line 20
\draw (i2) -| (part-M3.drain); % line 21
\draw (k2) -| (part-M3.source); % line 22
\draw (j1) -| (part-M3.gate); % line 23
\draw (i5) -| (part-M4.drain); % line 24
\draw (k5) -| (part-M4.source); % line 25
\draw (j4) -| (part-M4.gate); % line 26
\path (1.2,-4.193) rectangle (4.248,-3.007); % line 28
\node[anchor=west, inner sep=0, circuitnotemark, font=\large] at (1.2,-3.6) {X}; % line 28
\path (4.8,-4.193) rectangle (7.848,-3.007); % line 29
\node[anchor=west, inner sep=0, circuitnotemark, font=\large] at (4.8,-3.6) {X}; % line 29
\path (1.2,-8.993) rectangle (4.502,-7.807); % line 30
\node[anchor=west, inner sep=0, circuitnotemark, font=\large] at (1.2,-8.4) {X}; % line 30
\path (4.8,-8.993) rectangle (8.102,-7.807); % line 31
\node[anchor=west, inner sep=0, circuitnotemark, font=\large] at (4.8,-8.4) {X}; % line 31
\path (1.2,-13.793) rectangle (4.502,-12.607); % line 32
\node[anchor=west, inner sep=0, circuitnotemark, font=\large] at (1.2,-13.2) {X}; % line 32
\path (4.8,-13.793) rectangle (8.102,-12.607); % line 33
\node[anchor=west, inner sep=0, circuitnotemark, font=\large] at (4.8,-13.2) {X}; % line 33
\end{circuitikz}
\end{document}
